% Generated draft: requires mathematical and source review.
\documentclass[11pt]{article}
\usepackage[margin=25mm]{geometry}
\usepackage{fontspec}
\setmainfont{Latin Modern Roman}
\usepackage{amsmath,amssymb,mathtools,unicode-math}
\setmathfont{Latin Modern Math}
\usepackage{enumitem,needspace,xurl,hyperref}
\setlength{\parindent}{0pt}
\setlength{\parskip}{5pt}
\setlength{\emergencystretch}{3em}
\newcommand{\N}{\mathbb{N}}
\newcommand{\Z}{\mathbb{Z}}
\newcommand{\Q}{\mathbb{Q}}
\newcommand{\R}{\mathbb{R}}
\newcommand{\C}{\mathbb{C}}
\newcommand{\F}{\mathbb{F}}
\newcommand{\A}{\mathbb{A}}
\newcommand{\PP}{\mathbb{P}}
\newcommand{\E}{\mathbb{E}}
\newcommand{\eps}{\varepsilon}
\newcommand{\dd}{\,\mathrm{d}}
\begin{document}
\section*{UnsolvedMath ID: 30006071}
\textbf{Source problem code:} catalogue entry
\subsection{Definitions and mathematical statement}
Let \(n\geq 1\) be an integer and let \(A\in M_n(\Q(z))\). Write \(\partial=\dd/\dd z\) and consider the differential system
\[
(\partial-A)\mathbf f=0.
\]
Here, interpret ``all solutions are algebraic'' in the usual complex local sense: at every point \(z_0\in\C\) where \(A\) has no pole, every local holomorphic solution germ \(\mathbf f=(f_1,\ldots,f_n)^\mathsf T\) has each component algebraic over \(\C(z)\). Explicitly, for each \(i\) there is a nonzero polynomial \(P_i(T)\in\C(z)[T]\) such that \(P_i(f_i)=0\) as a germ of a meromorphic function. This interpretation allows arbitrary complex initial values; the source does not explicitly specify the field over which solutions are to be algebraic.

To define reduction, choose \(B\in M_n(\Z[z])\) and a nonzero \(q\in\Z[z]\) with \(A=B/q\). For each prime \(p\) such that \(\overline q\neq 0\), put
\[
A_p=\frac{\overline B}{\overline q}\in M_n(\F_p(z)),\qquad
D_p:\F_p(z)^n\longrightarrow\F_p(z)^n,\qquad
D_p(\mathbf v)=\frac{\dd\mathbf v}{\dd z}-A_p\mathbf v.
\]
The power \(D_p^p\) means composition of \(D_p\) with itself \(p\) times. Since \((\dd/\dd z)^p=0\) on \(\F_p(z)\), this composition is \(\F_p(z)\)-linear and is the \(p\)-curvature evaluated on \(\dd/\dd z\). Thus the vanishing condition means that \(D_p^p\) is the zero endomorphism of \(\F_p(z)^n\), rather than merely vanishing on horizontal solutions.

\textbf{Conjectural statement (Grothendieck--Katz \(p\)-curvature conjecture).} For every such \(n\) and \(A\), is the following equivalence true?
\[
\begin{gathered}
\text{Every local holomorphic solution of }(\partial-A)\mathbf f=0
\text{ is algebraic over }\C(z)
\\[2pt]
\Longleftrightarrow
\\[2pt]
\exists\text{ a finite set of primes }S\quad
\forall\text{ primes }p\notin S:\quad
\overline q\neq 0\ \text{ and }\ D_p^p=0.
\end{gathered}
\]
The finite exceptional set includes primes where the chosen presentation cannot be reduced.

\textbf{Unresolved source qualification.} The supplied clean statement specifies this linear differential-system conjecture. The supplied original excerpt also contains an incomplete formulation of a stronger nonlinear converse to Eisenstein's theorem, with missing conditions and an undefined function \(\omega(p)\). That additional conjecture cannot be stated precisely from this record and is not reconstructed here.
\subsection{Short English statement}
The conjecture asks whether every solution of a rational linear differential system is algebraic exactly when its \(p\)-curvature vanishes for all but finitely many primes. Algebraicity is interpreted here for complex local solution germs over \(\C(z)\). The supplied record labels the conjecture open; its incomplete nonlinear Eisenstein-related formulation requires separate clarification.
\subsection{Sources}
\begin{itemize}
\item Supplied local UnsolvedMath registry record, permanent ID \(30006071\), source problem code \texttt{catalogue entry}, titled \emph{The Grothendieck--Katz p-Curvature Conjecture}, in the Oberwolfach Reports Open Problems collection. The record lists proposed year 2024 and status \texttt{open}. Its \texttt{corrected\_verified} label and statement that the official report was checked are supplied provenance claims, not verification performed for this draft.
\item \emph{Anabelian Geometry and Representations of Fundamental Groups (2025)}, Oberwolfach Reports, DOI 10.4171/owr/2024/45, pp.~2621--2670, as cited in the supplied record: \url{https://doi.org/10.4171/owr/2024/45}. The literature-source label is ``Oberwolfach Report 45/2024.'' The supplied excerpt identifies Daniel Litt, jointly with Yeuk Hay Joshua Lam, in connection with ``On the converse to Eisenstein's last theorem''; this does not establish that they proposed the Grothendieck--Katz conjecture.
\item Lam--Litt, \emph{Algebraicity and integrality of solutions to differential equations}, supplied literature reference: \url{https://arxiv.org/abs/2501.13175}.
\item Literature assessment dated exactly \texttt{2026-08-22}: the record describes both the Grothendieck--Katz \(p\)-curvature conjecture and its stronger nonlinear Eisenstein-integrality conjecture as open in general, with substantial special classes proved. This is supplied dated curation derived from the OWR record, not an independently verified literature assessment. The record explicitly requests independent verification of time-sensitive status and cited claims.
\item The original excerpt attributes special-case results to Katz [3] for Picard--Fuchs equations, and to Chudnovsky--Chudnovsky [4], Bost [5], and Andr\'e [6] for solvable monodromy, and mentions other sporadic cases. It attributes an Eisenstein implication to [7]. These labels are reproduced from the excerpt; their corresponding bibliographic entries and precise hypotheses were not supplied.
\end{itemize}
\end{document}
